\documentclass[utf8]{frontiersSCNS} % for Science, Engineering and Humanities and Social Sciences articles

\usepackage{url,hyperref,lineno,microtype,subcaption}
\usepackage[onehalfspacing]{setspace}

\linenumbers


% Leave a blank line between paragraphs instead of using \\


\def\keyFont{\fontsize{8}{11}\helveticabold }
\def\firstAuthorLast{Sample {et~al.}} %use et al only if is more than 1 author
\def\Authors{Stefan Dirnstorfer}
% Affiliations should be keyed to the author's name with superscript numbers and be listed as follows: Laboratory, Institute, Department, Organization, City, State abbreviation (USA, Canada, Australia), and Country (without detailed address information such as city zip codes or street names).
% If one of the authors has a change of address, list the new address below the correspondence details using a superscript symbol and use the same symbol to indicate the author in the author list.
\def\Address{Laboratory X, Institute X, Department X, Organization X, City X , State XX (only USA, Canada and Australia), Country X }
% The Corresponding Author should be marked with an asterisk
% Provide the exact contact address (this time including street name and city zip code) and email of the corresponding author
\def\corrAuthor{Corresponding Author}

\def\corrEmail{stefan@thetaris.com}


\newcommand{\emoticon}[1]{\includegraphics[width=2cm]{gen/emoticon-#1}}


\begin{document}
\onecolumn
\firstpage{1}

\title[Running Title]{A parametric emoji for the representation of emotional states}

\author{\Authors} %This field will be automatically populated
\address{} %This field will be automatically populated
\correspondance{} %This field will be automatically populated

\extraAuth{}% If there are more than 1 corresponding author, comment this line and uncomment the next one.
%\extraAuth{corresponding Author2 \\ Laboratory X2, Institute X2, Department X2, Organization X2, Street X2, City X2 , State XX2 (only USA, Canada and Australia), Zip Code2, X2 Country X2, email2@uni2.edu}


\maketitle


\begin{abstract}

%%% Leave the Abstract empty if your article does not require one, please see the Summary Table for full details.
\section{}
For full guidelines regarding your manuscript please refer to \href{http://www.frontiersin.org/about/AuthorGuidelines}{Author Guidelines}.

As a primary goal, the abstract should render the general significance and conceptual advance of the work clearly accessible to a broad readership. References should not be cited in the abstract. Leave the Abstract empty if your article does not require one, please see \href{http://www.frontiersin.org/about/AuthorGuidelines#SummaryTable}{Summary Table} for details according to article type. 


\tiny
 \keyFont{ \section{Keywords:} keyword, keyword, keyword, keyword, keyword, keyword, keyword, keyword} %All article types: you may provide up to 8 keywords; at least 5 are mandatory.
\end{abstract}

\section{Introduction}

Many applications rely on the ability to represent a human emotion in computational form.

Current methods are either limited in the range of representable emotions or lack a clear interpretation of respective states.

The only widely established technique is to use facial expressions for human emotion

These can be subjectively and intuitively understood by everyone.
However, there is a lack of consensus on how to process such faces computationally.


Emoji are the most widely used system of facial expressions as they are now integrated in electronic devices and part of the unicode standard.
The extreme range of available expressions make them difficult to interpret.
Some emoji seem to be on a scale of increasing intensity, while others seem to be exact negations of each other. However, no clear consensus exists. Some emoji are even interpreted inconsistently accoss cultures


Many computational systems use a limited set of emotions when mapping a facial expression to emotional dimensions.
Ekman faces
Limited in scope. Pain, embarrassment, relief are missing.
Internal relationships are not faithfully reproduced.
Arousal, excitement and contempt are perfectly combinable and thus should not be an either or decision.
Other dimensions such as exceitement and anger are exclusive.
Henc, the values produced by such a representation system are not ver valuable.


Another system of facial expressions is the Wong-Baker FACES® Pain Rating Scale.
It's medical value and track record are beyond doubt and applied with wide success in pain treatment.
This system is absolutely clear in its interpretation.
The faces are ordered along an axis and can be turned into statistics, to eg measure an average emotion or its fluctuaton.
This feature make them a very powerful tool in pain therapy, as they satisfy the needs of experts and laymen alike.
However, the scale is only one-dimensional and the range of emotional states is limited to a single aspect of the perceivable space.


Maybe the earliest system of facial visualizations are the Chernoff faces.
These have a very clear numerical interpretation, as each movable part of a human face is controlled by a numeric value.
However, these faces lack an emotional interpretation. It would for example be difficult to define which faciescan be considered to perceive more pain. As such Chernoff faces could be shown to convey no additional information, when compared to other visualization techniques.
\cite{chernoff}

This paper suggests a facial visualization that includes 5 quantitative aspects of a human emotion.
The emotinal dimensions are derived from common psychological models.
The graphics definition is mostly derived from linear combinations of input coordinates.
The contribution of this representation is neither an improved model of emotional representation
nor is it a better set of emoticons.
What it does however contribute is compromise that might strike the right balance between psychological precision, emotional variety and graphical simplicity.

\section{The facial dimensions}

Only express emotional attributes
No physiological properties: age, gender, race
Inspired by established emotional dimensions
Graphical simplicity with linear displacements
Linearity is, of course, not a strict requirement, but in the absence of an objective scale a welcome guidance

\subsection{Valence and arousal}
Valence and arousal \cite{circumplex}

Linearly displaces all control points
Additionally combinable with all subsequent dimensions

\begin{figure}[ht!]
\begin{center}
\begin{tabular}{ccc}
\emoticon{0,100,50,0,50} &
\emoticon{50,100,50,0,50} &
\emoticon{100,100,50,0,50} \\
\emoticon{0,50,50,0,50} &
\emoticon{50,50,50,0,50} &
\emoticon{100,50,50,0,50} \\
\emoticon{0,0,50,0,50} &
\emoticon{50,0,50,0,50} &
\emoticon{100,0,50,0,50}
\end{tabular}
\end{center}
\caption{Facial expressions created with 3 different values for valence and arousal.}
\end{figure}


\subsection{Potency}

\begin{tabular}{cc|cc|cc}
\multicolumn{2}{c}{potency = 0.0} &
\multicolumn{2}{c}{potency = 0.5} &
\multicolumn{2}{c}{potency = 1.0} \\
\hline
\emoticon{46,22,0,0,50} &
\emoticon{39,50,0,0,50} &
\emoticon{62,42,50,0,50} &
\emoticon{39,30,50,0,50} &
\emoticon{92,43,100,0,50} &
\emoticon{90,2,100,0,50}
\\
\emoticon{3,84,0,0,50} &
\emoticon{97,12,0,0,50} &
\emoticon{18,100,50,0,50} &
\emoticon{92,27,50,0,50} &
\emoticon{8,10,100,0,50} &
\emoticon{66,94,100,0,50}
\end{tabular}

\subsection{Contempt}

\begin{tabular}{cc|cc|cc}
\multicolumn{2}{c}{contempt = 0.0} &
\multicolumn{2}{c}{contempt = 0.5} &
\multicolumn{2}{c}{contempt = 1.0} \\
\hline
\emoticon{44,21,27,0,50} &
\emoticon{14,55,67,0,50} &
\emoticon{86,19,88,50,50} &
\emoticon{91,38,25,50,50} &
\emoticon{80,14,29,100,50} &
\emoticon{11,92,41,100,50}
\\
\emoticon{95,43,98,0,50} &
\emoticon{93,57,59,0,50} &
\emoticon{27,42,38,50,50} &
\emoticon{71,13,7,50,50} &
\emoticon{10,10,90,100,50} &
\emoticon{72,27,82,100,50}
\end{tabular}

\subsection{Expression}

\section{Code}
The graphics of the presented facial expression is defined in SVG and Javascript computer language that can be understood by any modern web browser.

The coordinates and angles all derive linearly from input parameters defined above.

Hence, every coordinate of the face is defined by 6 numbers defining the default position and a displacement that is proportional to the input parameters.

This simplicity allows for a rational interpretation of the face, as each emotional dimension has a clear and independent geometric effect. Yet, the emotional state conveyed by each of the faces creates surprising new qualities as those features are combined.

\section{Conclusion}
There are many applications where emotional states need to be captured, processed or displayed.

Smaller emotional scope than emoji, as there exist some emojis for very particular emotions that do not fit in a dimensional model.
The facial expression has fewer dimensions than Chernoff faces which contain dimensions for, eg. size of ears that cannot be varied by people at will.
The dimensions of the presented face are inspired by scientific models, but do they do lack the depth and psychological justification of the eg. the Wong-Baker scale.
However, it is a middle ground between alternative systems.
Thus, it fills a gap that is left open by existing methods.
\cite{chernoff}

\section{References}

\subsection{WongBaker}
\url{https://en.wikipedia.org/wiki/Wong%E2%80%93Baker_Faces_Pain_Rating_Scale}

\subsection{DVPRS}

\url{https://academic.oup.com/painmedicine/article/16/11/2152/2460410}

\subsection{Integrated review}

\url{https://www.frontiersin.org/articles/10.3389/fpsyg.2016.02061/full}

\url{https://www.frontiersin.org/articles/10.3389/fpsyg.2019.02221/full}

\subsection{Noto emoji}

\url{https://www.google.com/get/noto/help/emoji/}

\url{https://www.researchgate.net/publication/266508208_What_the_Face_Reveals_Basic_and_Applied_Studies_of_Spontaneous_Expression_Using_the_Facial_Action_Coding_System_FACS}

\section*{Supplemental Data}
 \href{http://home.frontiersin.org/about/author-guidelines#SupplementaryMaterial}{Supplementary Material} should be uploaded separately on submission, if there are Supplementary Figures, please include the caption in the same file as the figure. LaTeX Supplementary Material templates can be found in the Frontiers LaTeX folder.

\section*{Data Availability Statement}
The datasets [GENERATED/ANALYZED] for this study can be found in the [NAME OF REPOSITORY] [LINK].
% Please see the availability of data guidelines for more information, at https://www.frontiersin.org/about/author-guidelines#AvailabilityofData

\bibliographystyle{frontiersinSCNS_ENG_HUMS}
\bibliography{faces}
 
\end{document}
